% !TEX root = ../main.tex
\section{World Access and the Limits of Training Exposure}
\label{sec:world-access}

Condition (2) of Definition~\ref{def:grounding} may look like an unnecessary
strengthening. If a claim appears in the training corpus, is that not already
grounding enough? It is not, and the reason is worth stating separately, because
the assumption that it is underwrites a good deal of confident practice.

Knowledge of contingent facts depends on an appropriate relation to evidence. A
model contains statistical traces of reports about the world, but training
exposure is not present evidentiary access. Writing $\TrainedOn(M, E)$ for the
relation that $M$ was trained on $E$, the following entailments all fail:
\begin{align}
  \TrainedOn(M, E) &\nentails \Authentic(E), \label{eq:not-authentic}\\
  \TrainedOn(M, E) &\nentails \Current(E, \tau), \label{eq:not-current}\\
  \TrainedOn(M, E) &\nentails \Applicable(E, j, u), \label{eq:not-applicable}\\
  \TrainedOn(M, E) &\nentails \Supports(E, \phi, c), \label{eq:not-supports}\\
  \TrainedOn(M, E) &\nentails \Traceable(E). \label{eq:not-traceable}
\end{align}

The first three are the substance of Section~\ref{sec:misapplied} restated at the
level of the corpus: a corpus is a snapshot, and a snapshot has a date, a
provenance, and a scope. The fourth is the point of
Section~\ref{sec:grounding-operator}: containing text that mentions $\phi$ is not
the same as containing evidence that favours $\phi$ for the decision at hand.

The fifth deserves emphasis, because it is where the engineering consequence
bites hardest. Training data influences parameter values. Parameter influence is
not an audit trail. It does not identify the particular source, edition,
jurisdiction, or effective date on which a given output should be relied
upon---and in a regulated setting, that identification is not a nice-to-have but
the thing being asked for. When a decision is questioned two years later, the
question is not whether the model was trained on good data. It is which document,
in which version, as of which date, supported this specific claim. A model that
has absorbed excellent sources and cannot say which one it used has not supplied
the answer.

This is the second and independent route to the conclusion of
Section~\ref{sec:measurement}. There, the argument was that the hallucination
rate is estimable only where the evidential base is explicit at inference time.
Here, the argument is that reconstructability requires the same thing, for
reasons that have nothing to do with measurement and everything to do with
accountability. Both converge on an architecture in which the evidence supporting
a claim is identified, versioned, and resolved to a context at the point of
assertion. Section~\ref{sec:warrant} sketches what that architecture has to
deliver; developing it is the business of a companion paper.
